% =====================================================================
%  Chapter 1 — The Investor and the Machine
% =====================================================================

\chapter{The Investor and the Machine}\idx{Investor}
\label{ch:investor}

This book is about one investor managing one portfolio\idx{Portfolio} with money that
belongs to that investor. Not a fund. Not an endowment with a quarterly
mark-to-model and a staff of twenty. Not a research division with a mandate
to publish. One person, one account, one set of decisions, taken over
decades, with a broker on the other end of them.

That constraint is the whole design. It is what makes the problem tractable
and what makes the answer unsatisfying, and a book that pretended otherwise
would be about a different problem. Let us be precise about what it buys and
what it costs.

\section{What one investor buys}

A single investor\idx{Single investor} faces no capacity constraint in any sense that matters.
The literature on implementation shortfall and on the cost of exploiting
anomalies is dominated by institutional results: roughly a trillion dollars
of live data across nineteen markets from 1998 to 2011, actual trading costs
found to be less than a tenth of prior estimates, and a clear ordering of
which anomalies survive the costs --- value and momentum scale, size does
not, and short-horizon reversals are the most capacity-constrained of all.
Those are the numbers of a pension fund trying to move a hundred million
dollars without moving the price against itself. They are the wrong numbers
for this reader.

The reader of this book is trying to move perhaps fifty thousand dollars in a
fund whose daily volume is measured in hundreds of millions, and who can
trade in the first fifteen minutes of the session without moving a single
basis point. In such a regime the marginal cost of a trade is the
bid--ask spread and the commission, and nothing else. The economically
significant frictions for this reader are not market impact at all. They are
the tax code, the bid--ask spread, and the behavioural failure modes of the
person doing the trading.

This is an unusually favourable position, and it is worth being clear about
how unusual. It means the book's central question is not \emph{can this
signal be traded?} but \emph{can this signal be traded by one person, without
a hedge fund?} That is a much narrower question, and it has a much better
answer.

\section{What one investor lacks}

Three things.

The first is \emph{diversification of judgement}. A fund with a research
staff has several independent opinions about the same question, and when they
disagree the disagreement is information. A single investor has one opinion,
and the only way to check it is to look at the data --- which is exactly what
this book spends two hundred pages on doing, and what the machine-learning\idx{Machine learning}
layer is for.

The second is \emph{the ability to separate signal from noise by
construction}. When a professional's backtest fails, it is usually because
they tried enough variants that one of them worked. The variants they tried
were not recorded, so the failure is invisible. The single investor who
follows this book is asked to do the opposite: to record every variant tried,
to declare the family before the test, and to be prepared to publish the
result even when it is that the machine adds nothing.

The third is \emph{time}. Not the reader's patience --- this book is not
asking anyone to watch a screen --- but the horizon over which a decision
must pay for itself. A single investor in the accumulation phase has a
twenty- to forty-year horizon and a fixed annual saving. A decision that
reduces the probability of ruin over that horizon is worth more than a
decision that raises the expected arithmetic return by a fraction of a
percent, because the arithmetic return is paid once, at the end, and the
ruin is paid whenever it happens.

\begin{takeaway}[The three asymmetries]
\begin{itemize}
\item \textbf{Cost.} For one investor, capacity is not a constraint.
Bid--ask spread and commission are the whole of the friction.
\item \textbf{Data.} For one investor, overfitting is the dominant risk.
There is no second opinion to catch it.
\item \textbf{Horizon.} For one investor with a long horizon, the variance
of the \emph{terminal} wealth, not the mean, is the thing to minimise.
\end{itemize}
\end{takeaway}

\section{The question, stated precisely}

The owner's question, in the form this book answers it, is:

\begin{quote}
\textit{Is there a machine-learning procedure that a single investor can
implement, that makes money and protects principal\idx{Principal} at the same time, and
that can be evaluated honestly enough that the investor knows which of those
two it is actually doing?}
\end{quote}

Three clauses, and each one is a trap.

\emph{Makes money} is ambiguous in a way that matters. It could mean the
arithmetic mean return exceeds the market's, which is what an academic
paper claims and what almost no live system delivers after costs. It could
mean the geometric mean exceeds the benchmark's, which is the only version
of the claim that a saver actually experiences, and which is strictly harder
than the arithmetic version, because volatility is a cost. This book takes
the geometric version. Every headline number in it is a compound annual
growth rate, and no claim is made in terms of the arithmetic mean except
where the arithmetic mean is the object of study.

\emph{Protects principal} is ambiguous in a way that is worse, because it
sounds like a claim that can be made and is usually made. Principal
protection is not a property a portfolio has; it is a property of a
distribution of outcomes, and every portfolio has one. What can be said
is that the portfolio has a bounded worst case, a drawdown the investor
can survive, and a recovery time the investor can wait out. Those are three
quantities, they are computable, and they are the ones this book computes.

\emph{Evaluated honestly} is the clause that decides whether the other two
mean anything. A machine-learning strategy can be evaluated honestly only
if the family of strategies tried is declared in advance, the data is free
of look-ahead, the costs and taxes are charged, and the selection effect is
subtracted from the result. Most published claims fail one or more of these
four tests, and a large fraction fail all four.

\section{Why the machine is the instrument and not the protagonist}

There is a version of this book in which the machine does the work and the
result is a handsome equity curve. That version is not available, for a
reason that is now well established in the literature and that this book
takes as its starting premise rather than as a controversy to be won.

The replication evidence is not encouraging for the ambitious version.
Jensen, Kelly and Pedersen, writing in the \emph{Journal of Finance} in
2023, surveyed a large literature of asset-pricing tests, reimplemented
hundreds of them, and found that a majority of published factors do
replicate out of sample --- and that the failures that do occur are
concentrated in a set of characteristics where replication is
methodologically hard. They also found that the returns of factors that
\emph{do} replicate have weakened over the past decade, and they made a
sharper methodological point: the standard errors used to evaluate a new
factor are too small, so that factors at the margin are being accepted at
rates well above the nominal five per cent.

McLean and Pontiff's result is the one that should govern a single
investor's expectations. They took 97 variables with published predictive
power, formed portfolios out of them, and then traded those portfolios
after publication rather than before it. Returns fell by 26 per cent
out of sample and by 58 per cent after publication, and 32 per cent of the
decline was attributable to the trading itself --- to the costs and
capacity of exploiting a signal that other people now know about.

And Goyal, Welch and Zafirov, surveying 2022, found that more than a third
of the predictors in the literature were no longer significant even
in the original sample.

\begin{commonmyths}
\item \textbf{Myth: ``Machine learning finds signals that humans miss.''}
Fact: what machine learning reliably finds is the subset of human-found
signals that survive a large number of trials. The finding is about
\emph{selection}, not about foresight.
\item \textbf{Myth: ``If the backtest works, the strategy works.''}
Fact: the backtest is the only place the strategy was ever tested. Every
published backtest that was later traded has underperformed its own
backtest; the median degradation reported in the literature is 26 per cent
and the post-publication degradation is 58 per cent.
\item \textbf{Myth: ``A more sophisticated model is a better model.''}
Fact: a more sophisticated model has more ways to be wrong and more
parameters to overfit. The sophistication that pays is the one that
prices its own uncertainty, and Chapter 20 shows what that looks like.
\end{commonmyths}

\section{What this book therefore does}

It builds one system, completely, and then takes it apart in front of the
reader.

Part I states the question and the evidence that bounds what any answer
could look like. Part II sets out the six facts about long-horizon investing
that any system must respect: the size of the premium, the sequence of
returns, myopic loss aversion, the behavioural traps, the frictions, and
the capacity position of a single investor. Part III builds the data layer
from raw price series, with the live window of every fund recorded and the
look-ahead traps named. Part IV builds the model layer from its objective
function upwards --- no library, no black box --- and then builds the
evaluation machinery that will be used to judge it. Part V builds the
portfolio layer: a stated strategic policy that the machine is not allowed
to touch, a bounded tilt inside each sleeve, and a risk governor. Part VI
assembles and runs the whole thing, with the tax engine attached, and reports
the result. Part VII reports what the result is worth.

The result is not the one the owner asked for. The machine-learning layer
adds very little to the strategic allocation and the risk governor, and the
selection tests cannot distinguish the system's return from the benchmark's.
The book says so, in the chapters where that is established, with the
arithmetic shown. What replaces the promise of alpha\idx{Alpha} is something smaller
and more reliable: a procedure that produces a sensible portfolio, knows how
little it knows, and shrinks its positions when the evidence is thin.

\begin{keytakeaways}
\item One investor, one portfolio, own money: capacity is not the binding
constraint, overfitting is.
\item ``Makes money'' means a geometric return after costs, taxes and
selection; ``protects principal'' means a bounded drawdown, a survivable
recovery time, and a stated worst twelve-month outcome.
\item The literature that would license a confident machine-learning alpha
does not exist: replication is partial, decay is documented, and a third of
published predictors have already died in sample.
\item The book therefore builds a system whose first job is to know what it
does not know.
\end{keytakeaways}
